\section{Step 40 Decide whether another block is needed}\label{chacha-s40}

Subtract the completed fragment size from the remaining count. A zero remainder takes the normal restore path immediately. With the 65-byte recipe the remainder is one, so execution continues toward the counter advance.

\subsection{Write and inspect}
Replace the body of region \texttt{stream\_remaining}. Retain its \texttt{\# BEGIN} and \texttt{\# END} comments and the supplied labels outside them.

\Needspace{8\baselineskip}
\MARSCodeLabel{MIPS assembly}
\begin{minted}[samepage=true]{mars-idea}
    dsubu $s2, $s2, $s3
    beqz $s2, stream_done
    nop
\end{minted}

Use recipe \texttt{s40-1} with the cumulative program through this edit. Inspect its named checkpoint before continuing to the temporary or normal return.

\subsection{Expected observation}
The first fragment is complete and \texttt{s2} is one; no second block has been generated yet.

\Needspace{9\baselineskip}
\subsection{Check your understanding}
Should a 64-byte request generate a second block?

\begin{itemize}
  \item \textbf{A} Yes; every loop needs a terminating empty block.
  \item \textbf{B} No; its remainder is zero.
  \item \textbf{C} Only when the first ciphertext byte is zero.
\end{itemize}
